
\chapter{Introduction}\label{sec-introduction}

This book is intended for a broad audience, ranging from novices to experts in data science and marketing. It serves as a bridge between the two fields, introducing foundational concepts and guiding readers through hands-on examples in R and RStudio, with a focus on applying text analytics to marketing contexts. At its heart is language, and language matters; it is a flow of data: When taking public transportation and hearing people chattering, this overlapping of sounds creates an almost indistinct noise, but within this is information being sent between travelers, and at times we can cut through the ''noise'' to decipher these messages. language, however, is especially durable when it is written down and leaves traces. When users express their opinions in writing on the internet they generate a huge stream of unstructured text, known as user-generated content (UGC). This mass of data is both rich and overwhelming. Selecting, collecting, and analyzing UGC is daunting and requires specific tools. In this book, we will show you how to do this by combining conceptual insights with practical instruction, beginning with the basic terminology using R programming techniques before progressing to real-world text analysis. The goal is to equip readers with the skills needed to tackle similar challenges when dealing with real data.

Obviously, the process of carrying out text analytics requires certain qualities which are universal for anyone who learns: curiosity, patience, and determination. The goal is to extract meaningful insights from several textual sources. This book focuses on practical applications such as analyzing online reviews, extracting key trends, and using text data to inform marketing decisions. But, given the right dataset, the principles described here can be applied to any scientific field which uses huge volumes of text. The text analyst's work is similar to that of a detective: finding enough evidence to capture insightful information. This work is also a bridge (another one) between disciplines such as linguistics, computer science, mathematics, and statistics.

Marketing is a field in a constant state of change, and in recent years it has been revolutionized by the ability to analyze text data on a large scale. Textual data can be represented by customer reviews, social media posts, opinions on forums, chats, or survey responses. Understanding these data points can reveal crucial information regarding the relationships between companies or brands and customers, through customer sentiment, preferences, and emerging trends. In this book, the focus on the methods of text analytics is tailored to marketing needs, providing examples and insights drawn from real-world scenarios.

Text mining is important in marketing for several reasons. Firstly, it helps us to help understand customer feedback, and marketers can use it to analyze customer reviews, feedback, and social media mentions to discern customers' sentiments and preferences. This information can then be used to inform product development, marketing campaigns, and customer service strategies. Secondly, an important point is to identify trends and insights. Marketers can analyze large amounts of unstructured data to identify emerging trends, pain points, and areas of opportunity. These can inform marketing strategies and help companies stay ahead of their competition. Thirdly, personalization is another crucial element that can help marketers segment customers and personalize marketing messages based on their preferences and behavior. And, last but not least, text mining can automate tasks like sentiment analysis, keyword extraction, and topic modeling, saving marketers time and resources.

In contemporary society and the digital economy, textual data has emerged as an important source of information for marketers, complementing traditional metrics in both richness and contextual relevance. As consumers increasingly engage with brands through digital platforms by writing online reviews, posting on social media, chatting with customer-service bots, or composing emails, the volume and variety of textual content have expanded at an unprecedented scale. This proliferation of unstructured language-based data presents both a challenge and an opportunity for marketing analysts: It is a challenge in terms of methodological complexity, and an opportunity to gain more nuanced context-sensitive insights into consumer-firm interactions, brand perceptions, and market dynamics.

Unlike structured numerical data, text captures the subtleties of human expression. Emotions, intentions, and nuanced preferences are often inaccessible through standard surveys or transactional records. Advances in natural language processing (NLP) and machine learning (ML) allow for scalable analyses of such text, enabling marketers to use a systematic, data-driven interpretation of language to support their decision-making. Consequently, textual analysis is becoming central to applications ranging from sentiment analysis and brand positioning, to customer segmentation and content marketing development.

A growing body of literature has explored the role of text analytics in marketing, and this book is closely aligned with several other contributions in this field. For example, Berger et al. (2020) provide a broad overview of how automated textual analysis can generate insights across various marketing domains, such as analyzing word of mouth, predicting brand crises, understanding consumer-firm interactions, and business-to-business communication. Similarly, Humphreys (2021) explores the use of automated text analysis in the context of consumer research and its potential for generating insights into marketing. There are several books which are oriented towards practice and use R for quantitative text-analysis applications. An example is Ledolter and VanderVelde (2021) who provide an accessible introduction to the quantitative analysis of textual data. Yildirim and Kübler (2023) target more advanced marketing students and practitioners, while Ashlee (2016) covers core marketing topics for social media.

Our book is located in the same space but aims to provide a more applied treatment of the automated text analysis techniques introduced by Humphreys and Wang (2018). While their work presents an excellent overview of approaches such as dictionary-based methods, sentiment analysis, and topic modeling, we provide extensive R code and practical examples to implement these techniques. Like Ledolter and VanderVelde (2021), we adopt a practical, hands-on approach to textual analysis, but with a clear focus on marketing-specific applications. It also complements Ashlee's (2016) work on social media by diving deeper into quantitative text-analysis methodologies for marketing applications.

\section{The Tools Used}\label{used-tools}

This book uses the programming language R and the RStudio integrated development environment (IDE) extensively as tools for text analytics, and whether you are new to R or a more experienced user, this book offers helpful practical guidance. R provides a rich ecosystem of packages for manipulating, analyzing, and visualizing textual data. In the first few chapters, we cover essential built-in commands, but as we proceed, we will introduce libraries like \texttt{tidyverse} and \texttt{ggplot} to conduct our analyses. Before we start, however, we need to install these tools on our computers, and at a conceptual level, this is a universal task, independent of your Operating System. Very briefly, R can be installed from \url{https://cran.r-project.org/} for all the main OS, either in precompiled binary form (the best choice for the vast majority of people) or from source code (for the bold and daring). The next step is to install RStudio from \url{https://posit.co}; this is the graphical interface that uses R to enhance its potential. Although initially developed to support only R, it now also supports Python and is designed to accommodate other programming languages as well. We suggest that you use the open-source edition from the Posit website. Be aware that the installation operation must be done in a specific sequence: first R and then RStudio.

\subsection{Packages}\label{packages}

Packages are a system of extensions that allow us to extend the basic capabilities of R by enhancing data analysis, graphical representation, or the ease of use. Once installed inside the R environment on your computer, the functions existing in the packages can be employed by simply using the package's name in the function \texttt{library(nameofthepackage)}. The difference between a package and a library is like getting a book from a bookstore: You buy some books, and they are wrapped by the clerk in a package, you bring them home, unwrap them and put the books in your home library. This means you do not need to go and fetch the books (the functions of a package) all the time from the bookstore (the internet) and bring them home (download them), you can simply read them at home (load the library from your computer).

Depending on whether you use the command line or the RStudio interface, there are different ways to install packages in R. Using a command line, you have three ways to install them:

\begin{itemize}
\tightlist
\item
  From the CRAN network (R Core Team 2025; the recommended procedure)

  \begin{itemize}
  \tightlist
  \item
    \texttt{install.packages("packagename")}
  \end{itemize}
\item
  From Github (GitHub, Inc. 2025) as a repository (some packages are not present on the CRAN network but are still useful)

  \begin{itemize}
  \tightlist
  \item
    \texttt{install.packages("devtools")}
  \item
    \texttt{install\_github("username/packagename")}
  \end{itemize}
\item
  From a file downloaded from the internet

  \begin{itemize}
  \tightlist
  \item
    \texttt{install.packages("path\_and\_source",\ repos\ =\ NULL,\ type="source")}
  \item
    \texttt{install.packages("\textasciitilde{}/Downloads/packagex.zip",\ repos=NULL,\ type="source")}
  \item
    \texttt{install.packages(file.choose(),\ repos=NULL)\ \#\ on\ Windows}
  \end{itemize}
\end{itemize}

If you employ packages in your code that have not been
installed, RStudio will prompt you whenever it sees that the library
corresponding to a package is present, but not the package itself.

R does not automatically check for new packages, so check from time
to time that you are up-to-date, especially if you installed the packages a while ago. Be aware that elements might not operate better
after an update. Some procedures or functions might be modified or even
removed, so we advise exercising caution when updating packages. Here
are some simple functions that are useful when dealing with the packages
already installed on your computer.

\begin{Shaded}
\begin{Highlighting}[]
\CommentTok{\# Check for installed packages}
\FunctionTok{installed.packages}\NormalTok{()}

\CommentTok{\# List all packages, this also checks if an update is available}
\FunctionTok{old.packages}\NormalTok{()}

\CommentTok{\# Update all (available) packages}
\FunctionTok{update.packages}\NormalTok{()}

\CommentTok{\# Update, without prompts for permission/clarification}
\FunctionTok{update.packages}\NormalTok{(}\AttributeTok{ask =} \ConstantTok{FALSE}\NormalTok{)}

\CommentTok{\# Updating only one package is like overwriting on top of the old one, it uses the same command}
\FunctionTok{install.packages}\NormalTok{(}\StringTok{"package"}\NormalTok{)}
\end{Highlighting}
\end{Shaded}

Two of the newer packages which give R users a lot more power to deal with and combine functions in R, are \texttt{magrittr} and \texttt{tidyverse}. The former, uses a set of ``pipes'' represented by the \texttt{\%\textgreater{}\%} operator to use the output of one function as the input of the next. In this small example, we introduce you to the difference this package can make in writing cleaner and more compact code. If the task is to take a vector of numbers, perform the square root of them, round the results to two decimal places, and keep the numbers to greater than 4, using a traditional approach, you would write:

\begin{Shaded}
\begin{Highlighting}[]
\NormalTok{numbers }\OtherTok{\textless{}{-}} \FunctionTok{c}\NormalTok{(}\DecValTok{1}\NormalTok{, }\DecValTok{4}\NormalTok{, }\DecValTok{9}\NormalTok{, }\DecValTok{16}\NormalTok{, }\DecValTok{25}\NormalTok{)}

\CommentTok{\# Square root}
\NormalTok{sqrt\_numbers }\OtherTok{\textless{}{-}} \FunctionTok{sqrt}\NormalTok{(numbers)}

\CommentTok{\# Round to two decimals}
\NormalTok{rounded\_numbers }\OtherTok{\textless{}{-}} \FunctionTok{round}\NormalTok{(sqrt\_numbers, }\DecValTok{2}\NormalTok{)}

\CommentTok{\# Filter numbers greater than 4}
\NormalTok{filtered\_numbers }\OtherTok{\textless{}{-}}\NormalTok{ rounded\_numbers[rounded\_numbers }\SpecialCharTok{\textgreater{}} \DecValTok{2}\NormalTok{]}

\FunctionTok{print}\NormalTok{(filtered\_numbers)}
\end{Highlighting}
\end{Shaded}

\begin{verbatim}
[1] 3 4 5
\end{verbatim}

The same result can be achieved with \texttt{magrittr} like this, using fewer variables and so less code:

\begin{Shaded}
\begin{Highlighting}[]
\FunctionTok{library}\NormalTok{(magrittr)}

\NormalTok{numbers }\OtherTok{\textless{}{-}} \FunctionTok{c}\NormalTok{(}\DecValTok{1}\NormalTok{, }\DecValTok{4}\NormalTok{, }\DecValTok{9}\NormalTok{, }\DecValTok{16}\NormalTok{, }\DecValTok{25}\NormalTok{)}

\NormalTok{filtered\_numbers }\OtherTok{\textless{}{-}}\NormalTok{ numbers }\SpecialCharTok{\%\textgreater{}\%}
  \FunctionTok{sqrt}\NormalTok{() }\SpecialCharTok{\%\textgreater{}\%}               \CommentTok{\# Step 1: Take the square root}
  \FunctionTok{round}\NormalTok{(}\DecValTok{2}\NormalTok{) }\SpecialCharTok{\%\textgreater{}\%}             \CommentTok{\# Step 2: Round to two decimals}
\NormalTok{  .[. }\SpecialCharTok{\textgreater{}} \DecValTok{2}\NormalTok{]                 }\CommentTok{\# Step 3: Filter numbers greater than 4}

\FunctionTok{print}\NormalTok{(filtered\_numbers)}
\end{Highlighting}
\end{Shaded}

\begin{verbatim}
[1] 3 4 5
\end{verbatim}

In contrast, \texttt{tidyverse} uses \texttt{magrittr} heavily and is a collection of packages to be used in data manipulation, exploration, and visualization. For example, some of the key packages, among many others, are \texttt{dplyr} which is used to filter, summarize, or join, \texttt{ggplot2} to visualize data, \texttt{tidyr} is used to reshape datasets, and \texttt{readr} to import CSV files.

Now that you have some knowledge of how to install R and packages, let us focus on how to obtain the accompanying data for this book.

\section{Dataset Overview and Companion
Website}\label{dataset-overview-and-companion-website}

Collecting data is a time-consuming activity because data are unstructured and ubiquitous on the internet. But these data must be funneled into a structured dataset so that our lives as analysts become much easier. Figure~\ref{fig-overview} shows how data has to be transformed into a more manageable format before it can be analyzed, such as CSV, JSON, a tabular format, or any other structured form.

\begin{figure}

\centering{

\includegraphics[width=0.9\textwidth,height=\textheight]{figure/chapter1/overview.png}

}

\caption{\label{fig-overview}Data from several sources, unstructured to
structured \newline Authors' own figure}

\end{figure}%

For the dataset in this book in the main we decided to use one single source, an older collection of Amazon user reviews of tablets (Xin Wang, Mai, and Chiang 2014). This will serve as the context for applying the text analytics techniques presented here. A companion website – \href{https://danieldanlan.github.io/}{https://danieldanlab.github.io/} – was added as an essential resource for anyone interested in using this book as an instruction manual. It hosts the dataset and other related files, additional examples of code and scripts, exercises for hands-on practice, and updates and supplementary material to extend your learning. Readers are encouraged to visit the website frequently to access these materials and enhance their understanding of the concepts discussed in the book. The goal is to equip the reader with constantly updated information, keeping the book as the foundation for learning and the website for updates, a problem-solving environment, suggestions for code, and any other helpful information which may arise.

\section{Structure of the Book}\label{structure-of-the-book}

The remainder of this book is structured to guide the reader through the fundamentals and advanced techniques of text analytics in the following way:

Chapter 2 (\emph{Textual Data}) deals with the basics of working with text in R, including commands, structures, and preprocessing. Basic R commands and regular expressions are presented with examples. Novice users learn how to clean text, deal with common problems in UGC, and prepare data for analysis. This chapter is the foundation for dealing with text data and prepares the ground for advanced analytical techniques.

Chapter 3 (\emph{Obtaining Textual Data}) tackles strategies for obtaining data for analysis, from simply importing text files and web scraping, to working with application program interfaces (APIs) and other commonly used methods. By the end of the chapter, readers will understand where to look for relevant data and how to bring them into R for preprocessing and analysis.

Chapter 4 (\emph{Basic Text Analysis}) starts by introducing techniques for text representation (e.g., a bag-of-words or TF-IDF) and visualization. Basic operations such as word frequencies, how to normalize text, and how to perform basic statistical analyses are introduced here. This chapter serves as an intermediary step between data preparation and more advanced analytical procedures, giving readers a first view of the potential of the extraction of information from text processing.

Chapter 5 (\emph{Clustering}) analyzes methods for grouping text or documents, including hierarchical and k-means clustering, as well as density-based approaches. Meaningful patterns can be found by understanding how to group similar documents or words. Using these algorithms, we demonstrate how clustering reveals structures within preprocessed documents. The chapter also introduces a more direct link between text analytics and marketing applications.

Chapter 6 (\emph{Classification}) deals with supervised learning models for categorizing text data. Regular machine learning procedures, such as k-nearest neighbors, support vector machines, ensemble learning, as well as simple neural-network models that are used in other fields can also be applied to text. The chapter underlines how these methods help to categorize data and connect classification to marketing applications.

Chapter 7 (\emph{Topic Modeling}) explores an essential element in many marketing applications, helping the reader to identify latent themes in text data using methods such as latent Dirichlet allocation. Readers can learn how to extract latent topics in data and understand their importance for marketing strategies. Visualization techniques are also discussed.

Chapter 8 (\emph{Sentiment Analysis}) deals with analyzing sentiment over time and understanding customer emotions, extracting sentiment polarization. It also introduces aspect-based sentiment analysis concepts, how to build sentiment word networks, and how to detect negative customer feedback. By the end of the chapter, the reader will have grasped the value of sentiment analysis for understanding customer perceptions.

Chapter 9 (\emph{Named Entity Recognition and Summarization}) deals with extracting key pieces of information, such as names, brands, and locations. Summarizing text is also tackled. The examples show how these methods are used in text analytics to enhance efficiency and understanding, without the need to explore the entire dataset.

The last chapter (\emph{Advanced Techniques}) introduces the concept of word embedding, explores cosine similarity as a metric for comparing textual data, and transformer models. Methods such as these which bring artificial intelligence (AI) into analytics allow readers to explore semantic relationships in documents. This chapter is the link that connects traditional natural language processing (NLP) methods with advances in AI through neural networks and text vectorization.

\section{Key Features and Tools
Covered}\label{key-features-and-tools-covered}

The book offers a practical, hands-on approach to text analytics. Key highlights include real-world examples using marketing data, code snippets, step-by-step tutorials, the integration of basic and advanced R packages, as well as free access to a dedicated website for extended learning. NLP is a vast and fast-changing domain; it comprises subsections such as natural language understanding (NLU) and natural language generation (NLG), both of which are integral to analyzing and producing text. This book focuses on practical NLP applications, particularly in marketing contexts, but the techniques can be transferred to other domains, such as psychology or journalism.

\section{Challenges in Text
Analytics}\label{challenges-in-text-analytics}

Analyzing text data presents unique challenges, which we discuss in the book. For example, dealing with unstructured and inconsistent data formats is one of them, addressing encoding issues across operating systems (Windows, Linux or MacOS) and across different languages. Managing linguistic variations, such as British vs. American English, is another key challenge. The book offers strategies and tools for navigating these complexities effectively.

\section{Setting Expectations for
Readers}\label{setting-expectations-for-readers}

This book is positioned at the intersection between computer science and marketing analytics, placing textual data at the core of insight generation and decision making. The target audiences are data and computer scientists with an interest in marketing applications, as well as marketing analysts seeking to improve their computational abilities and gain access to a robust set of tools for processing textual data. The objective is to equip researchers and practitioners with the techniques and frameworks needed to extract value from text data systematically. Throughout the book, we illustrate the applications primarily using a dataset containing customer feedback provided in online reviews. No prior programming experience is required; the book starts with fundamental concepts and progresses to advanced topics. Readers are encouraged to explore the exercises and access the companion website to maximize their learning experience. For the novice user, a good free resource to learn R is ``An Introduction to R'', which can be found here: \url{https://intro2r.com/}
